\documentclass[10pt,a4paper]{amsart}
\usepackage{amsmath,amssymb,amsthm,graphicx}
% Retained preamble for the Stacks Project chapter {CH} (shared project preamble).
% The chapter file pinned-original-source.tex only carries \input{{preamble}} on its first line,
% so no preamble text lives inside the chapter itself.  The Stacks Project uses one single shared
% preamble for every chapter: the snapshot stages it three times (sources/stacks-chow/preamble.tex,
% sources/stacks-etale-cohomology/preamble.tex and focused-seeds/AGP/sources/stacks-snapshot-20261011/
% preamble.tex) and those copies are identical in size.  The block below is copied mechanically,
% line for line, from sources/stacks-chow/preamble.tex lines 14-18 and 154-217: the Xy-pic setup and
% the native theorem declarations together with the mathematical macros.  Nothing was added, removed
% or redefined; the wrapper class and the amsmath/amssymb/amsthm/graphicx base packages are supplied
% by the assembler.
% For commutative diagrams we use Xy-pic
% Chapter: stacks-spaces-more-morphisms.tex -- More on Morphisms of Spaces.
\usepackage[all]{xy}

% We use 2cell for 2-commutative diagrams.
\xyoption{2cell}

% Theorem environments.
%
\theoremstyle{plain}
\newtheorem{theorem}[subsection]{Theorem}
\newtheorem{proposition}[subsection]{Proposition}
\newtheorem{lemma}[subsection]{Lemma}

\theoremstyle{definition}
\newtheorem{definition}[subsection]{Definition}
\newtheorem{example}[subsection]{Example}
\newtheorem{exercise}[subsection]{Exercise}
\newtheorem{situation}[subsection]{Situation}

\theoremstyle{remark}
\newtheorem{remark}[subsection]{Remark}
\newtheorem{remarks}[subsection]{Remarks}

\numberwithin{equation}{subsection}

% Macros
%
\def\lim{\mathop{\mathrm{lim}}\nolimits}
\def\colim{\mathop{\mathrm{colim}}\nolimits}
\def\Spec{\mathop{\mathrm{Spec}}}
\def\Hom{\mathop{\mathrm{Hom}}\nolimits}
\def\Ext{\mathop{\mathrm{Ext}}\nolimits}
\def\SheafHom{\mathop{\mathcal{H}\!\mathit{om}}\nolimits}
\def\SheafExt{\mathop{\mathcal{E}\!\mathit{xt}}\nolimits}
\def\Sch{\mathit{Sch}}
\def\Mor{\mathop{\mathrm{Mor}}\nolimits}
\def\Ob{\mathop{\mathrm{Ob}}\nolimits}
\def\Sh{\mathop{\mathit{Sh}}\nolimits}
\def\NL{\mathop{N\!L}\nolimits}
\def\CH{\mathop{\mathrm{CH}}\nolimits}
\def\proetale{{pro\text{-}\acute{e}tale}}
\def\etale{{\acute{e}tale}}
\def\QCoh{\mathit{QCoh}}
\def\Ker{\mathop{\mathrm{Ker}}}
\def\Im{\mathop{\mathrm{Im}}}
\def\Coker{\mathop{\mathrm{Coker}}}
\def\Coim{\mathop{\mathrm{Coim}}}

% Boxtimes
%
\DeclareMathSymbol{\boxtimes}{\mathbin}{AMSa}{"02}

%
% Macros for moduli stacks/spaces
%
\def\QCohstack{\mathcal{QC}\!\mathit{oh}}
\def\Cohstack{\mathcal{C}\!\mathit{oh}}
\def\Spacesstack{\mathcal{S}\!\mathit{paces}}
\def\Quotfunctor{\mathrm{Quot}}
\def\Hilbfunctor{\mathrm{Hilb}}
\def\Curvesstack{\mathcal{C}\!\mathit{urves}}
\def\Polarizedstack{\mathcal{P}\!\mathit{olarized}}
\def\Complexesstack{\mathcal{C}\!\mathit{omplexes}}
% \Pic is the operator that assigns to X its picard group, usage \Pic(X)
% \Picardstack_{X/B} denotes the Picard stack of X over B
% \Picardfunctor_{X/B} denotes the Picard functor of X over B
\def\Pic{\mathop{\mathrm{Pic}}\nolimits}
\def\Picardstack{\mathcal{P}\!\mathit{ic}}
\def\Picardfunctor{\mathrm{Pic}}
\def\Deformationcategory{\mathcal{D}\!\mathit{ef}}
\usepackage{cleveref}
\crefname{originallemma}{original lemma}{original lemmas}
\Crefname{originallemma}{Original lemma}{Original lemmas}
\makeatletter
\newlabel{lemma-quasi-finite-separated-pass-through-finite-addendum}{{0874}{source}{}{sourcepointer-1}{}}
\newlabel{lemma-quasi-finite-separated-pass-through-finite-addendum@cref}{{[originallemma][0][]0874}{[1][0][]source}{}{}{}}
\makeatother
\crefname{originalreference}{original reference}{original references}
\Crefname{originalreference}{Original reference}{Original references}
\makeatletter
\newlabel{more-algebra-lemma-fitting-ideal-basics}{{07ZA}{source}{}{sourcepointer-2}{}}
\newlabel{more-algebra-lemma-fitting-ideal-basics@cref}{{[originalreference][0][]07ZA}{[1][0][]source}{}{}{}}
\makeatother
\crefname{originalreference}{original reference}{original references}
\Crefname{originalreference}{Original reference}{Original references}
\makeatletter
\newlabel{spaces-descent-lemma-finite-finitely-presented-module}{{0610}{source}{}{sourcepointer-3}{}}
\newlabel{spaces-descent-lemma-finite-finitely-presented-module@cref}{{[originalreference][0][]0610}{[1][0][]source}{}{}{}}
\makeatother
\crefname{originalreference}{original reference}{original references}
\Crefname{originalreference}{Original reference}{Original references}
\makeatletter
\newlabel{spaces-limits-lemma-extend}{{0853}{source}{}{sourcepointer-4}{}}
\newlabel{spaces-limits-lemma-extend@cref}{{[originalreference][0][]0853}{[1][0][]source}{}{}{}}
\makeatother
\crefname{originalreference}{original reference}{original references}
\Crefname{originalreference}{Original reference}{Original references}
\makeatletter
\newlabel{spaces-morphisms-lemma-closed-immersion-finite-presentation}{{084Q}{source}{}{sourcepointer-5}{}}
\newlabel{spaces-morphisms-lemma-closed-immersion-finite-presentation@cref}{{[originalreference][0][]084Q}{[1][0][]source}{}{}{}}
\makeatother
\crefname{originalreference}{original reference}{original references}
\Crefname{originalreference}{Original reference}{Original references}
\makeatletter
\newlabel{spaces-morphisms-lemma-composition-finite-presentation}{{03XQ}{source}{}{sourcepointer-6}{}}
\newlabel{spaces-morphisms-lemma-composition-finite-presentation@cref}{{[originalreference][0][]03XQ}{[1][0][]source}{}{}{}}
\makeatother
\crefname{originalreference}{original reference}{original references}
\Crefname{originalreference}{Original reference}{Original references}
\makeatletter
\newlabel{spaces-morphisms-lemma-composition-integral}{{03ZR}{source}{}{sourcepointer-7}{}}
\newlabel{spaces-morphisms-lemma-composition-integral@cref}{{[originalreference][0][]03ZR}{[1][0][]source}{}{}{}}
\makeatother
\title{Pushing a finite type ideal along a quasi-compact separated etale morphism}
\author{The Stacks Project Authors}
\date{stacks-project:spaces-more-morphisms.tex:local-snapshot-20261011; GNU FDL 1.2 or later}
\begin{document}
\maketitle
\noindent\textit{Editorial note.} Counted span is the author's own proof environment, source lines 8715-8775, immediately following the statement at 8697-8713.\par\medskip
\noindent\textit{Editorial note.} The Fitting ideal computation is displayed with all three equalities inside the span.\par\medskip
\noindent\textit{Editorial note.} Six cross-chapter pointers (More on Algebra, Descent on Spaces, Limits of Spaces, Morphisms of Spaces) are recorded in assembly.external\_references.\par\medskip
\noindent\textit{Allowed standard prerequisites (editorial).}\begin{itemize}
\item Fitting ideals and their behaviour under finite locally free morphisms. More on Algebra, lemma more-algebra-lemma-fitting-ideal-basics, cited at proof lines 8770-8772. Computes the Fitting ideal of the pushed module and is the algebraic core of the construction.
\item Extension of finite modules over thickenings. Limits of Spaces, lemma spaces-limits-lemma-extend, cited in the proof. Extends the relevant finite module from a closed piece to a neighbourhood, which is how the ideal I' is produced.
\item Finiteness properties of closed immersions and compositions. Morphisms of Spaces lemmas on finite presentation of closed immersions, compositions and integral morphisms, cited in the proof. Keeps the constructed ideal of finite type along the auxiliary closed subscheme Z and its finite piece Z'.
\item The pass-through-finite addendum. chapter lemma lemma-quasi-finite-separated-pass-through-finite-addendum, cited in the proof. Supplies the finite piece Z' of the auxiliary construction.
\end{itemize}\medskip
\ifdefined\hypertarget\hypertarget{sourcepointer-1}{}\fi
\noindent\textit{Source reference.} Original lemma 0874: sources/stacks-spaces-more-morphisms/pinned-original-source.tex, source lines 7806--7806. Navigation pointer only: the counted proof cites this statement by \textbackslash{}ref as a standing input; the deduction the counted proof has to make is carried out inside the counted span, and no part of that deduction is delegated to this pointer.\par\medskip
\ifdefined\hypertarget\hypertarget{sourcepointer-2}{}\fi
\noindent\textit{Source reference.} Original reference 07ZA: focused-seeds/geometry/S9-spaces-push-ideal/ref-tags.txt, source lines 4--4. Navigation pointer only: the counted proof cites this statement by \textbackslash{}ref as a standing input; the deduction the counted proof has to make is carried out inside the counted span, and no part of that deduction is delegated to this pointer.\par\medskip
\ifdefined\hypertarget\hypertarget{sourcepointer-3}{}\fi
\noindent\textit{Source reference.} Original reference 0610: focused-seeds/geometry/S9-spaces-push-ideal/ref-tags.txt, source lines 5--5. Navigation pointer only: the counted proof cites this statement by \textbackslash{}ref as a standing input; the deduction the counted proof has to make is carried out inside the counted span, and no part of that deduction is delegated to this pointer.\par\medskip
\ifdefined\hypertarget\hypertarget{sourcepointer-4}{}\fi
\noindent\textit{Source reference.} Original reference 0853: focused-seeds/geometry/S9-spaces-push-ideal/ref-tags.txt, source lines 6--6. Navigation pointer only: the counted proof cites this statement by \textbackslash{}ref as a standing input; the deduction the counted proof has to make is carried out inside the counted span, and no part of that deduction is delegated to this pointer.\par\medskip
\ifdefined\hypertarget\hypertarget{sourcepointer-5}{}\fi
\noindent\textit{Source reference.} Original reference 084Q: focused-seeds/geometry/S9-spaces-push-ideal/ref-tags.txt, source lines 7--7. Navigation pointer only: the counted proof cites this statement by \textbackslash{}ref as a standing input; the deduction the counted proof has to make is carried out inside the counted span, and no part of that deduction is delegated to this pointer.\par\medskip
\ifdefined\hypertarget\hypertarget{sourcepointer-6}{}\fi
\noindent\textit{Source reference.} Original reference 03XQ: focused-seeds/geometry/S9-spaces-push-ideal/ref-tags.txt, source lines 8--8. Navigation pointer only: the counted proof cites this statement by \textbackslash{}ref as a standing input; the deduction the counted proof has to make is carried out inside the counted span, and no part of that deduction is delegated to this pointer.\par\medskip
\ifdefined\hypertarget\hypertarget{sourcepointer-7}{}\fi
\noindent\textit{Source reference.} Original reference 03ZR: focused-seeds/geometry/S9-spaces-push-ideal/ref-tags.txt, source lines 9--9. Navigation pointer only: the counted proof cites this statement by \textbackslash{}ref as a standing input; the deduction the counted proof has to make is carried out inside the counted span, and no part of that deduction is delegated to this pointer.\par\medskip
\begin{lemma}
\label{lemma-push-ideal}
Let $S$ be a scheme. Let $X$ be a quasi-compact and quasi-separated
algebraic space over $S$. Let $\varphi : W \to X$ be a quasi-compact
separated \'etale morphism. Let $U \subset X$ be a quasi-compact open
subspace. Let $\mathcal{I} \subset \mathcal{O}_W$ be a finite type
quasi-coherent sheaf of ideals such that
$V(\mathcal{I}) \cap \varphi^{-1}(U) = \emptyset$.
Then there exists a finite type quasi-coherent sheaf of ideals
$\mathcal{J} \subset \mathcal{O}_X$ such that
\begin{enumerate}
\item $V(\mathcal{J}) \cap U = \emptyset$, and
\item $\varphi^{-1}(\mathcal{J})\mathcal{O}_W = \mathcal{I} \mathcal{I}'$
for some finite type quasi-coherent ideal
$\mathcal{I}' \subset \mathcal{O}_W$.
\end{enumerate}
\end{lemma}

\begin{proof}
Choose a factorization $W \to Y \to X$ where $j : W \to Y$ is a
quasi-compact open immersion and $\pi : Y \to X$ is a finite
morphism of finite presentation
(Lemma \ref{lemma-quasi-finite-separated-pass-through-finite-addendum}).
Let $V = j(W) \cup \pi^{-1}(U) \subset Y$. Note that
$\mathcal{I}$ on $W \cong j(W)$ and $\mathcal{O}_{\pi^{-1}(U)}$
glue to a finite type quasi-coherent sheaf of ideals
$\mathcal{I}_1 \subset \mathcal{O}_V$. By
Limits of Spaces, Lemma \ref{spaces-limits-lemma-extend}
there exists a finite type quasi-coherent sheaf of ideals
$\mathcal{I}_2 \subset \mathcal{O}_Y$ such that
$\mathcal{I}_2|_V = \mathcal{I}_1$. In other words,
$\mathcal{I}_2 \subset \mathcal{O}_Y$
is a finite type quasi-coherent sheaf of ideals such that
$V(\mathcal{I}_2)$ is disjoint from $\pi^{-1}(U)$ and
$j^{-1}\mathcal{I}_2 = \mathcal{I}$. Denote $i : Z \to Y$
the corresponding closed immersion which is of finite presentation
(Morphisms of Spaces, Lemma
\ref{spaces-morphisms-lemma-closed-immersion-finite-presentation}).
In particular the composition $\tau = \pi \circ i : Z \to X$ is finite
and of finite presentation
(Morphisms of Spaces, Lemmas
\ref{spaces-morphisms-lemma-composition-finite-presentation} and
\ref{spaces-morphisms-lemma-composition-integral}).

\medskip\noindent
Let $\mathcal{F} = \tau_*\mathcal{O}_Z$ which we think of as
a quasi-coherent $\mathcal{O}_X$-module. By
Descent on Spaces, Lemma
\ref{spaces-descent-lemma-finite-finitely-presented-module}
we see that $\mathcal{F}$ is a finitely presented $\mathcal{O}_X$-module.
Let $\mathcal{J} = \text{Fit}_0(\mathcal{F})$. (Insert reference to
fitting modules on ringed topoi here.) This is a finite type quasi-coherent
sheaf of ideals on $X$ (as $\mathcal{F}$ is of finite presentation, see
More on Algebra, Lemma \ref{more-algebra-lemma-fitting-ideal-basics}).
Part (1) of the lemma holds because $|\tau|(|Z|) \cap |U| = \emptyset$
by our choice of $\mathcal{I}_2$ and because the $0$th Fitting ideal
of the trivial module equals the structure sheaf. To prove (2) note that
$\varphi^{-1}(\mathcal{J})\mathcal{O}_W = \text{Fit}_0(\varphi^*\mathcal{F})$
because taking Fitting ideals commutes with base change.
On the other hand, as $\varphi : W \to X$ is separated and \'etale
we see that $(1, j) : W \to W \times_X Y$ is an open and closed immersion.
Hence $W \times_Y Z = V(\mathcal{I}) \amalg Z'$ for some finite and
finitely presented morphism of algebraic spaces $\tau' : Z' \to W$.
Thus we see that
\begin{align*}
\text{Fit}_0(\varphi^*\mathcal{F}) & =
\text{Fit}_0((W \times_Y Z \to W)_*\mathcal{O}_{W \times_Y Z}) \\
& =
\text{Fit}_0(\mathcal{O}_W/\mathcal{I}) \cdot
\text{Fit}_0(\tau'_*\mathcal{O}_{Z'}) \\
& =
\mathcal{I} \cdot \text{Fit}_0(\tau'_*\mathcal{O}_{Z'})
\end{align*}
the second equality by
More on Algebra, Lemma \ref{more-algebra-lemma-fitting-ideal-basics}
translated in sheaves on ringed topoi.
Setting $\mathcal{I}' = \text{Fit}_0(\tau'_*\mathcal{O}_{Z'})$
finishes the proof of the lemma.
\end{proof}

\end{document}
